\section{Background}
\label{sec:background}

\subsection{Knowledge Distillation}
\label{sec:bg-kd}

Knowledge distillation transfers the behaviour of a trained teacher network to a
smaller student by training the student to match the teacher's softened output
distribution rather than the hard labels alone \cite{hinton2015distilling}. The
softened targets, obtained by raising the softmax temperature, carry the relative
confidence the teacher assigns across classes, a richer signal than a one-hot
label, and the student minimizes a weighted sum of the ordinary cross-entropy and
a distillation term that aligns the two distributions. A compact student can thus
reach an accuracy it would not attain from the labels alone, which is what makes
distillation attractive when the model to be trained is small.

The benefit, however, is not monotone in the teacher's strength. When the teacher
is much larger than the student, the student cannot reproduce its distribution and
accuracy degrades instead of improving, an effect attributed to the capacity gap
between the two networks \cite{cho2019efficacy}. The gap can be bridged with
intermediate models of decreasing size: a teaching assistant is distilled from
the large teacher and in turn teaches the student, so that each step spans a
smaller capacity gap \cite{mirzadeh2020takd}. This staged distillation is the
principle behind the teacher-assistant scheme described in
Section~\ref{sec:method-kd}.

\subsection{Neural Network Quantization}
\label{sec:bg-quant}

Quantization replaces the floating-point weights and activations of a network
with low-precision fixed-point values, reducing both the memory the model occupies
and the cost of its arithmetic. A quantizer maps a real value to a finite integer
grid defined by a scale, a zero point when the representation is asymmetric, and a
bit-width. The smaller the grid, the cheaper the hardware, but also the larger the
rounding and clipping error that the network must tolerate.

Two strategies are common. Post-training quantization starts from an already
trained floating-point network and fixes the quantizer parameters afterwards. In
practice, this means calibrating activation ranges on representative data and
mapping each weight tensor to the nearest values on the low-bit grid, by scaling,
rounding to integer codes, and clipping to the representable range. It is cheap
because it does not require another optimization run, but at low bit-widths the
induced error is often too large for the original weights to absorb.
Quantization-aware training instead inserts fake quantizers during training. A
model can still be warm-started from a floating-point checkpoint, but the
low-precision operations are simulated in the forward pass while full-precision
parameters are retained for the gradient update. The network therefore learns
weights that tolerate the precision it will use at inference
\cite{jacob2018quantization}. It costs more to train but loses far less accuracy,
particularly at low bit-widths, so we use it throughout, through Brevitas, a
PyTorch library that inserts quantized layers and exports the trained model to a
quantized ONNX graph \cite{pappalardo2023brevitas}.

\subsection{FPGA Dataflow Inference}
\label{sec:bg-fpga}

FINN implements a quantized network as a streaming dataflow accelerator. Each
layer becomes its own hardware block, and the blocks form a pipeline through which
activations flow, so that every layer is busy on a different piece of data at the
same time \cite{umuroglu2017finn}. This organization suits low-precision networks,
because their arithmetic reduces to small fixed-point operations that map directly
onto the fabric's lookup tables and DSP slices, and the fewer the bits, the less
logic each operation costs. The original FINN targeted binarized networks; its
successor FINN-R generalizes the flow to arbitrary low bit-widths, and is the
toolchain we use \cite{blott2018finnr}.

The generated graph contains both compute blocks and stream-management blocks.
The compute blocks perform the layer arithmetic, mainly through matrix-vector
activation units for convolutional and fully connected layers and thresholding
blocks for quantized activations. The stream-management blocks prepare and move
data between them: padding and input-generation blocks arrange feature maps for
convolutions, data-width converters adapt how many bits are carried per cycle,
and FIFOs buffer values between neighbouring blocks with different rates. The
important point is that the non-arithmetic blocks are not free. In a fully
on-chip accelerator, weights, thresholds, and inter-layer buffers all compete for
FPGA memory, so stream management can become as relevant as the arithmetic itself.

The accelerator exposes a few parameters that trade throughput against resource
use. For matrix-vector activation units, folding is controlled by the PE and SIMD
factors. PE replicates output-channel processing, while SIMD replicates
input-channel processing inside each dot product. Raising either factor reduces
the number of cycles needed by the corresponding layer, but consumes more lookup
tables, registers, DSP slices, and often more on-chip memory for weights or
buffers. Between consecutive blocks, FIFO buffers absorb the mismatch in their
production and consumption rates and keep the pipeline from stalling. Because the
accelerator keeps all weights and these FIFOs on chip, the limited on-chip block
RAM bounds what fits, together with the logic; which of the two binds first
depends on the network and the folding chosen.
